%% ============================================================================
%% HEADER BLOCK
%% ============================================================================
%% Appendix:     PM
%% Category:     REF (Reference/How-To)
%% Name:         PARAMETH
%% Title:        Parameter Methodology Reference
%% Version:      1.0
%% Date:         2026-01-23
% Language: English
%% Dependencies: EST, CAL, WHERE, JJJ, CMP, Chapter 4x
%%
%% PURPOSE: Single entry point for ALL parameterization methodology in EBF.
%% This appendix consolidates and navigates the distributed parameterization
%% documentation - it does NOT repeat content but provides the complete map.
%% ============================================================================
%
% CROSS-REFERENCE MAP:
%   → EST (CORE-WHERE): Epistemics, Three-Tier Hierarchy, BBB Axioms
%   → CAL (METHOD-LLMMC): LLM Monte Carlo Pipeline
%   → WHERE: Parameter Repository
%   → JJJ (METHOD-PSG): Practical Sourcing Guide
%   → CMP (METHOD-COMP): Complementarity Estimation
%   → Chapter 4x: Calibration vs. Simulation Philosophy
%   → Chapters 17-20: Intervention Parameterization
%
% CHAPTER LINKAGE:
%   Primary: Chapter 4 (Empirical Foundations)
%   Secondary: Chapter 8 (Mathematical), Chapters 17-20 (Interventions)

\chapter{PM REF-PARAMETH: Parameter Methodology Reference}
\label{app:ref-pm}

%% ============================================================================
%% HEADER BOX
%% ============================================================================

\begin{tcolorbox}[colback=white, colframe=black!50, title=Appendix Information]
\begin{tabular}{ll}
\textbf{Appendix:} & PM \\
\textbf{Category:} & REF (Reference/How-To) \\
\textbf{Name:} & PARAMETH \\
\textbf{Version:} & 1.0 \\
\textbf{Primary Chapter:} & 4 (Empirical Foundations) \\
\textbf{Dependencies:} & EST, CAL, WHERE, JJJ, CMP \\
\end{tabular}
\end{tcolorbox}

%% ============================================================================
%% ABSTRACT
%% ============================================================================

\begin{abstract}
This appendix serves as the \textbf{single entry point} for understanding how to parameterize any EBF model component. Parameterization methodology is distributed across 10+ appendices and 6+ chapters - this reference consolidates the complete landscape into one navigable map. \textbf{Core principle:} EBF uses \textit{calibration on experimental evidence}, not simulation. Three tiers exist: (1) Hard Empirical, (2) Meta-Analytical, (3) Synthetic/LLMMC. This appendix answers: "I need a parameter value - where do I start and what method applies?"
\end{abstract}

%% ============================================================================
%% QUICK REFERENCE
%% ============================================================================

\begin{tcolorbox}[colback=blue!5!white, colframe=blue!75!black, title=Quick Reference: Parameter Methodology]
\small
\begin{tabular}{ll}
\textbf{Philosophy:} & Calibration, not Simulation (Chapter 4x) \\
\textbf{Hierarchy:} & Tier 1 $\succ$ Tier 2 $\succ$ Tier 3 (EST Axiom UNMAPPED_BBB-1) \\
\textbf{LLMMC Pipeline:} & Appendix UNMAPPED_CAL \\
\textbf{Literature Values:} & Appendix UNMAPPED_EST, WHERE \\
\textbf{Practical Guide:} & Appendix UNMAPPED_JJJ \\
\textbf{Complementarity $\gamma$:} & Appendix UNMAPPED_CMP \\
\textbf{Papers Database:} & \texttt{bcm\_master.bib} (1,922+ papers) \\
\end{tabular}
\end{tcolorbox}

%% ============================================================================
%% SCOPE BOX
%% ============================================================================

\begin{tcolorbox}[colback=gray!5!white, colframe=gray!75!black, title=Appendix Scope:]
\textbf{Ziel:} Provide a single navigation hub for EBF parameterization methodology.

\textbf{In-Scope:}
\begin{itemize}[nosep]
    \item Decision tree for method selection
    \item Parameter-type routing to specific appendices
    \item Three-tier hierarchy overview
    \item LLMMC pipeline summary
    \item Worked example demonstrating the full workflow
\end{itemize}

\textbf{Out-of-Scope (delegated to other appendices):}
\begin{itemize}[nosep]
    \item Complete tier hierarchy theory → Appendix UNMAPPED_EST
    \item Full LLMMC implementation details → Appendix UNMAPPED_CAL
    \item Parameter value tables → Appendix UNMAPPED_WHERE
    \item Complementarity estimation theory → Appendix UNMAPPED_CMP
\end{itemize}

\textbf{Lieferobjekte:} Navigation map, decision tree, parameter-type routing table, worked example.
\end{tcolorbox}

%% ============================================================================
%% FUNDAMENTAL QUESTION
%% ============================================================================

\begin{tcolorbox}[colback=orange!5!white, colframe=orange!75!black, title=Fundamental Question]
\textbf{How do I parameterize any EBF model component?}

Starting from "I need a value for parameter X", this appendix guides the practitioner through:
(1) which tier applies, (2) which method to use, (3) which appendix has the details.
\end{tcolorbox}

%% ============================================================================
%% MOTIVATION: WHY THIS APPENDIX?
%% ============================================================================

\section{Motivation: Why a Consolidation Appendix?}
\label{sec:pm-motivation}

\textbf{The Problem:} EBF's parameterization methodology is comprehensive but distributed:

\begin{itemize}[nosep]
    \item \textbf{EST} defines the epistemics and three-tier hierarchy
    \item \textbf{CAL} implements the LLMMC pipeline
    \item \textbf{WHERE} provides the parameter repository
    \item \textbf{JJJ} guides practical sourcing decisions
    \item \textbf{CMP} addresses complementarity estimation specifically
    \item \textbf{Chapter 4x} establishes the philosophical foundation
    \item \textbf{Chapters 17-20} apply parameterization to interventions
\end{itemize}

\textbf{The Gap:} No single document answers: "I need to parameterize $\gamma_{FS}$ for a health intervention in Germany - where do I start?"

\textbf{This Appendix Provides:}
\begin{enumerate}[nosep]
    \item Complete navigation map to all parameterization sources
    \item Decision tree for method selection
    \item Parameter-type routing (which appendix for which parameter)
    \item Integration overview showing how sources connect
\end{enumerate}

%% ============================================================================
%% CORE PRINCIPLE: CALIBRATION NOT SIMULATION
%% ============================================================================

\section{Core Principle: Calibration, Not Simulation}
\label{sec:pm-principle}

\begin{tcolorbox}[colback=red!5!white, colframe=red!75!black, title=Fundamental Distinction]
\textbf{EBF uses CALIBRATION on experimental evidence, NOT LLM-based simulation.}

\begin{itemize}[nosep]
    \item \textbf{Calibration:} Extract parameter values from published experimental results
    \item \textbf{Simulation:} Ask LLMs to role-play human decision-making
\end{itemize}

LLMs fail at behavioral simulation (see Chapter 4x for evidence). They succeed at synthesizing expert knowledge about \textit{what experiments have found}.
\end{tcolorbox}

\textbf{Reference:} Chapter 4x ``Calibration Not Simulation'' provides the complete philosophical foundation with empirical evidence for why this distinction is critical.

%% ============================================================================
%% THE THREE-TIER HIERARCHY
%% ============================================================================

\section{The Three-Tier Hierarchy}
\label{sec:pm-hierarchy}

\begin{tcolorbox}[colback=green!5!white, colframe=green!75!black, title=BBB Axiom UNMAPPED_BBB-1: Evidential Precedence]
$$\text{Tier 1} \succ \text{Tier 2} \succ \text{Tier 3}$$
Hard empirical data always supersedes meta-analytical estimates, which supersede synthetic priors.
\end{tcolorbox}

\subsection{Tier 1: Hard Empirical}

\begin{itemize}[nosep]
    \item \textbf{Source:} Direct experimental results (RCTs, field experiments)
    \item \textbf{Characteristics:}
        \begin{itemize}[nosep]
            \item Specific context (population, domain, intervention)
            \item Published effect sizes with confidence intervals
            \item Pre-registered designs preferred
        \end{itemize}
    \item \textbf{Example:} $\lambda = 2.25$ from Kahneman \& Tversky (1992)
    \item \textbf{Epistemic Tag:} \texttt{[EMP]}
    \item \textbf{Primary Source:} Appendix UNMAPPED_EST, \texttt{bcm\_master.bib}
\end{itemize}

\subsection{Tier 2: Meta-Analytical}

\begin{itemize}[nosep]
    \item \textbf{Source:} Systematic reviews, meta-analyses
    \item \textbf{Characteristics:}
        \begin{itemize}[nosep]
            \item Aggregated across studies
            \item Heterogeneity quantified ($I^2$, $\tau^2$)
            \item Domain-specific bounds available
        \end{itemize}
    \item \textbf{Example:} $\beta \in [0.7, 0.9]$ from hyperbolic discounting meta-analysis
    \item \textbf{Epistemic Tag:} \texttt{[THR]}
    \item \textbf{Primary Source:} Appendix UNMAPPED_EST, WHERE
\end{itemize}

\subsection{Tier 3: Synthetic Priors (Including LLMMC)}

\begin{itemize}[nosep]
    \item \textbf{Source:} Expert judgment, LLM Monte Carlo, analogical reasoning
    \item \textbf{Characteristics:}
        \begin{itemize}[nosep]
            \item No direct empirical measurement
            \item Based on theoretical considerations
            \item Higher uncertainty, requires sensitivity analysis
        \end{itemize}
    \item \textbf{Example:} $\gamma_{FS} = 0.35$ from LLMMC with 1000 experts
    \item \textbf{Epistemic Tag:} \texttt{[LLM]} or \texttt{[HYP]}
    \item \textbf{Primary Source:} Appendix UNMAPPED_CAL (LLMMC pipeline)
\end{itemize}

\textbf{Full Specification:} Appendix UNMAPPED_EST provides the complete axiomatic treatment of the three-tier hierarchy.

%% ============================================================================
%% MASTER DECISION TREE
%% ============================================================================

\section{Master Decision Tree: How to Parameterize}
\label{sec:pm-decision-tree}

\begin{tcolorbox}[colback=yellow!5!white, colframe=yellow!75!black, title=The Parameterization Decision Tree]
\begin{verbatim}
START: I need parameter value for X
       │
       ▼
┌──────────────────────────────────────────┐
│ STEP 1: Search bcm_master.bib            │
│ python scripts/search_bibliography.py X  │
└──────────────────────────────────────────┘
       │
       ├── Found RCT/Field Experiment? ──────► Tier 1 [EMP]
       │   → Use published effect size         (Done)
       │   → Reference: Appendix UNMAPPED_EST
       │
       ├── Found Meta-Analysis? ─────────────► Tier 2 [THR]
       │   → Use aggregated estimate           (Done)
       │   → Reference: Appendix UNMAPPED_EST, WHERE
       │
       └── Nothing directly applicable?
           │
           ▼
┌──────────────────────────────────────────┐
│ STEP 2: Check Parameter Repository       │
│ Appendix UNMAPPED_WHERE / EST tables              │
└──────────────────────────────────────────┘
       │
       ├── Found validated value? ───────────► Use with context
       │   → Check domain/population match     adjustment
       │   → Reference: Appendix UNMAPPED_WHERE
       │
       └── No match?
           │
           ▼
┌──────────────────────────────────────────┐
│ STEP 3: Apply LLMMC Pipeline             │
│ /r-score --demo or --full                │
│ Reference: Appendix UNMAPPED_CAL                  │
└──────────────────────────────────────────┘
       │
       ├── R-Score > 1.2? ───────────────────► GO [LLM]
       │   → Use LLMMC estimate
       │
       ├── 0.8 < R-Score < 1.2? ─────────────► PILOT
       │   → Collect additional data
       │
       └── R-Score < 0.8? ───────────────────► NO-GO
           → Parameter too uncertain
           → Redesign required
\end{verbatim}
\end{tcolorbox}

%% ============================================================================
%% PARAMETER-TYPE ROUTING
%% ============================================================================

\section{Parameter-Type Routing}
\label{sec:pm-routing}

Different parameters require different estimation approaches. Use this table to route to the correct appendix:

\begin{center}
\begin{tabular}{llll}
\toprule
\textbf{Parameter} & \textbf{Symbol} & \textbf{Primary Appendix} & \textbf{Method} \\
\midrule
Loss Aversion & $\lambda$ & EST, WHERE & Literature (Tier 1/2) \\
Present Bias & $\beta$ & EST, WHERE & Literature (Tier 1/2) \\
Discount Factor & $\delta$ & EST, WHERE & Literature (Tier 1) \\
Complementarity & $\gamma_{ij}$ & \textbf{CMP} & Structural estimation \\
Context Modulation & $\Psi$ & WHERE, WEN & Domain-specific \\
Phase Affinity & $\alpha$ & Ch.18, CAL & LLMMC + validation \\
Segment Multiplier & $\sigma$ & Ch.19, CAL & Subgroup analysis \\
FEPSDE Weights & $\omega_d$ & Ch.10, WAT & Domain calibration \\
Intervention Effect & $d$, $\Delta$ & CAL & LLMMC pipeline \\
R-Score & $R$ & \textbf{CAL} & LLMMC output \\
\bottomrule
\end{tabular}
\end{center}

\subsection{Special Case: Complementarity Coefficients ($\gamma$)}

Complementarity parameters require special treatment due to identification challenges:

\begin{itemize}[nosep]
    \item \textbf{Primary Reference:} Appendix UNMAPPED_CMP (METHOD-COMP)
    \item \textbf{Methods:}
        \begin{itemize}[nosep]
            \item Structural estimation with exclusion restrictions
            \item Instrumental variables for identification
            \item Revealed preference approaches
        \end{itemize}
    \item \textbf{Challenge:} Selection confounding between utility dimensions
    \item \textbf{Solution:} CMP provides identification strategies
\end{itemize}

%% ============================================================================
%% THE LLMMC PIPELINE (SUMMARY)
%% ============================================================================

\section{The LLMMC Pipeline (Summary)}
\label{sec:pm-llmmc}

When Tier 1/2 evidence is unavailable, the LLM Monte Carlo (LLMMC) pipeline provides Tier 3 estimates.

\begin{tcolorbox}[colback=blue!5!white, colframe=blue!75!black, title=LLMMC Pipeline Overview]
\textbf{Input:} Intervention specification (context, mechanism, population)

\textbf{Process:}
\begin{enumerate}[nosep]
    \item Generate $N$ expert perspectives (default: 1000)
    \item Each expert provides effect estimate on calibrated scale
    \item Apply scale-specific calibration formula
    \item Compute R-Score with uncertainty propagation
\end{enumerate}

\textbf{Output:} $R$-Score with GO/PILOT/NO-GO recommendation

\textbf{Full Pipeline:} Appendix UNMAPPED_CAL (METHOD-LLMMC)

\textbf{CLI:} \texttt{/r-score --demo} or \texttt{/r-score --full}
\end{tcolorbox}

\subsection{Calibration Formulas}

\textbf{For d/SMD Scale:}
$$\alpha = \frac{d_{\text{LLM}} \cdot \sigma_{\text{DV}}}{X_{\max} - X_{\min}} \cdot k_{\text{domain}}$$

\textbf{For Percentage-Point Scale:}
$$\alpha = \frac{\Delta_{\text{LLM}}}{100} \cdot \frac{1}{\bar{X}_{\text{baseline}}} \cdot k_{\text{domain}}$$

\textbf{Full derivation and validation:} Appendix UNMAPPED_CAL, Section 3.

%% ============================================================================
%% EPISTEMIC STATUS TAGS
%% ============================================================================

\section{Epistemic Status Tags}
\label{sec:pm-tags}

Every parameter in EBF must carry an epistemic status tag:

\begin{center}
\begin{tabular}{lll}
\toprule
\textbf{Tag} & \textbf{Meaning} & \textbf{Tier} \\
\midrule
\texttt{[EMP]} & Empirical - direct experimental evidence & Tier 1 \\
\texttt{[THR]} & Theoretical - meta-analysis or established theory & Tier 2 \\
\texttt{[LLM]} & LLM-derived via LLMMC pipeline & Tier 3 \\
\texttt{[ILL]} & Illustrative - pedagogical, not for decisions & N/A \\
\texttt{[HYP]} & Hypothetical - placeholder for sensitivity analysis & Tier 3 \\
\bottomrule
\end{tabular}
\end{center}

\textbf{Rule:} In any production model, parameters tagged \texttt{[ILL]} or \texttt{[HYP]} must be replaced with \texttt{[EMP]}, \texttt{[THR]}, or \texttt{[LLM]} values before deployment.

%% ============================================================================
%% COMPLETE SOURCE MAP
%% ============================================================================

\section{Complete Source Map}
\label{sec:pm-source-map}

\begin{center}
\begin{tabular}{lll}
\toprule
\textbf{Appendix} & \textbf{Category} & \textbf{Content} \\
\midrule
EST & CORE-WHERE & Epistemics, Three-Tier, BBB Axioms, 200+ values \\
WHERE & CORE/REF & Parameter Repository, context adjustment \\
CAL & METHOD-LLMMC & LLM Monte Carlo Pipeline, R-Score \\
JJJ & METHOD-PSG & Practical Sourcing Guide, worked examples \\
CMP & METHOD-COMP & Complementarity $\gamma$ estimation \\
BJ & METHOD-DOMAINVAL & Cross-cultural validation \\
EVD & METHOD-EVIDENCE & Evidence tier requirements \\
OPS & METHOD & Operationalization protocols \\
MCD & METHOD & Methodological foundations \\
\midrule
Ch. 4 & -- & Empirical Foundations overview \\
Ch. 4x & -- & Calibration vs. Simulation philosophy \\
Ch. 8 & -- & Mathematical framework \\
Ch. 17-20 & -- & Intervention parameterization \\
\bottomrule
\end{tabular}
\end{center}

%% ============================================================================
%% WORKED EXAMPLE
%% ============================================================================

\section{Worked Example: Parameterizing a Health Intervention}
\label{sec:pm-worked-example}

\textbf{Scenario:} Design a diabetes self-management intervention in Germany. Need parameters for:
\begin{itemize}[nosep]
    \item Loss aversion $\lambda$ for health outcomes
    \item Present bias $\beta$ for lifestyle changes
    \item Complementarity $\gamma_{PS}$ between Physical and Social dimensions
    \item Expected effect size for information provision
\end{itemize}

\subsection{Step 1: Search Literature}

\begin{verbatim}
python scripts/search_bibliography.py "loss aversion health"
python scripts/search_bibliography.py "present bias diabetes"
\end{verbatim}

\textbf{Result:} Found Tier 1 evidence for $\lambda = 2.1$ in health domain (specific to medical decisions).

\subsection{Step 2: Check Parameter Repository}

Consulting Appendix UNMAPPED_EST/WHERE tables:
\begin{itemize}[nosep]
    \item $\beta = 0.82$ [THR] - meta-analysis across health behaviors
    \item $\delta = 0.95$ [EMP] - German population specific
\end{itemize}

\subsection{Step 3: Complementarity via CMP}

For $\gamma_{PS}$ (Physical-Social complementarity):
\begin{itemize}[nosep]
    \item No direct Tier 1/2 evidence
    \item Apply CMP methodology: analogical reasoning from exercise studies
    \item Estimate: $\gamma_{PS} = 0.25$ [THR] with high uncertainty
\end{itemize}

\subsection{Step 4: LLMMC for Intervention Effect}

\begin{verbatim}
/r-score --intervention "diabetes education" --context "Germany, 50+"
\end{verbatim}

\textbf{Output:}
\begin{itemize}[nosep]
    \item $d_{\text{LLM}} = 0.35$ (moderate effect)
    \item $R = 1.45$ $\rightarrow$ GO
    \item Reference: 12 similar interventions in literature
\end{itemize}

\subsection{Final Parameter Set}

\begin{center}
\begin{tabular}{llll}
\toprule
\textbf{Parameter} & \textbf{Value} & \textbf{Tag} & \textbf{Source} \\
\midrule
$\lambda$ & 2.1 & [EMP] & Literature (Tier 1) \\
$\beta$ & 0.82 & [THR] & Meta-analysis (Tier 2) \\
$\delta$ & 0.95 & [EMP] & German study (Tier 1) \\
$\gamma_{PS}$ & 0.25 & [THR] & CMP analogy (Tier 2) \\
$d_{\text{info}}$ & 0.35 & [LLM] & LLMMC (Tier 3) \\
\bottomrule
\end{tabular}
\end{center}

%% ============================================================================
%% INTEGRATION WITH DEVELOPER PIPELINE
%% ============================================================================

\section{Integration with Developer Pipeline}
\label{sec:pm-integration}

This appendix integrates with the Developer Pipeline (Appendix UNMAPPED_DP) at Stage 3 (Calibration):

\begin{verbatim}
DP Stage 3: Calibration
├── Method A: BBB Literature → Use EST/WHERE
├── Method B: LLMMC → Use CAL Pipeline (this appendix routes here)
└── Method C: Empirical → Field data collection
\end{verbatim}

\textbf{Decision Matrix:}
\begin{itemize}[nosep]
    \item \textbf{New context, established intervention:} Method A (Literature)
    \item \textbf{Novel intervention, no prior data:} Method B (LLMMC)
    \item \textbf{Specific population, resources available:} Method C (Empirical)
    \item \textbf{Default:} A $\rightarrow$ B $\rightarrow$ C (cascading)
\end{itemize}

%% ============================================================================
%% KEY RESULTS
%% ============================================================================

\section{Key Results}
\label{sec:pm-results}

\begin{tcolorbox}[colback=green!5!white, colframe=green!75!black, title=Key Results]
\begin{enumerate}
    \item \textbf{Single Entry Point:} This appendix consolidates navigation to 10+ parameterization sources
    \item \textbf{Decision Tree:} Clear routing from "I need a parameter" to specific method
    \item \textbf{Parameter-Type Routing:} Each parameter type has a designated primary appendix
    \item \textbf{Tier Hierarchy:} Tier 1 $\succ$ Tier 2 $\succ$ Tier 3 enforced systematically
    \item \textbf{Epistemic Transparency:} All parameters tagged [EMP]/[THR]/[LLM]/[ILL]/[HYP]
\end{enumerate}
\end{tcolorbox}

%% ============================================================================
%% OPEN ISSUES
%% ============================================================================

\section{Open Issues}
\label{sec:pm-open-issues}

\begin{enumerate}
    \item \textbf{Tier Boundary Ambiguity:} Some estimates fall between tiers (e.g., small-sample RCT vs. large meta-analysis). Need clearer decision rules.

    \item \textbf{Context Transfer:} When is a Tier 1 estimate from Context A applicable to Context B? Currently handled ad-hoc in WHERE.

    \item \textbf{LLMMC Validation:} More ground-truth comparisons needed between LLMMC predictions and actual experimental results.

    \item \textbf{Parameter Updating:} Protocol for updating parameters when new evidence arrives. Currently manual.

    \item \textbf{Uncertainty Propagation:} How parameter uncertainty flows through to final predictions. Partially addressed in CAL.
\end{enumerate}

%% ============================================================================
%% CRITICAL FOUNDATIONS
%% ============================================================================

\section{Critical Foundations}
\label{sec:pm-foundations}

\begin{itemize}
    \item \textbf{UNMAPPED_BBB-1 (Evidential Precedence):} Tier 1 $\succ$ Tier 2 $\succ$ Tier 3 - Appendix UNMAPPED_EST
    \item \textbf{Calibration Principle:} LLMs calibrate on evidence, don't simulate behavior - Chapter 4x
    \item \textbf{Epistemic Transparency:} All parameters must be tagged - This appendix, Section \ref{sec:pm-tags}
    \item \textbf{Complementarity Identification:} $\gamma$ requires special methods - Appendix UNMAPPED_CMP
    \item \textbf{R-Score Threshold:} GO/PILOT/NO-GO at 1.2/0.8 - Appendix UNMAPPED_CAL
\end{itemize}

%% ============================================================================
%% SUMMARY
%% ============================================================================

\section{Summary}
\label{sec:pm-summary}

Appendix UNMAPPED_PM serves as the \textbf{navigation hub} for EBF's distributed parameterization methodology. Rather than repeating content, it provides:

\begin{enumerate}[nosep]
    \item The philosophical foundation (calibration, not simulation)
    \item The three-tier hierarchy with routing rules
    \item A decision tree for method selection
    \item Parameter-type specific guidance
    \item Complete cross-reference map to all sources
\end{enumerate}

\textbf{Start here when parameterizing.} Follow the decision tree in Section \ref{sec:pm-decision-tree}, then navigate to the specific appendix for detailed methodology.

%% ============================================================================
%% GLOSSARY
%% ============================================================================

\section{Glossary}
\label{sec:pm-glossary}

Key terms follow Appendix G (Glossary) master definitions. Key additions for this appendix:

\begin{description}[nosep]
    \item[Calibration:] Extracting parameter values from published experimental evidence
    \item[Simulation:] Asking models to role-play decision-making (NOT used in EBF)
    \item[LLMMC:] LLM Monte Carlo - method for Tier 3 estimates
    \item[R-Score:] Reliability score from LLMMC, thresholds: GO $>$ 1.2, PILOT $>$ 0.8
    \item[Epistemic Tag:] [EMP]/[THR]/[LLM]/[ILL]/[HYP] classification
\end{description}

%% ============================================================================
%% REFERENCES
%% ============================================================================

\section{References}
\label{sec:pm-references}

This appendix references the master bibliography (\texttt{bcm\_master.bib}).
\nocite{bcm_master}

\subsection{Meta-Analyses and Parameter Reviews (Tier 2 Sources)}

\begin{itemize}[nosep]
    \item \textbf{DellaVigna (2009)} - ``Psychology and Economics: Evidence from the Field'' \citep{dellavigna2009psychology} - \textit{Comprehensive meta-analysis of behavioral parameters across domains}
    \item \textbf{Frederick, Loewenstein \& O'Donoghue (2002)} - ``Time Discounting and Time Preference: A Critical Review'' \citep{frederick2002time} - \textit{Definitive review of discounting parameters [Tier 1]}
    \item \textbf{Cohen et al. (2020)} - ``Measuring Time Preferences'' \citep{cohen2020measuring} - \textit{Modern update on time preference estimation}
    \item \textbf{DellaVigna (2020)} - ``Structural Behavioral Economics'' \citep{dellavigna2020structural} - \textit{Handbook chapter on structural estimation}
\end{itemize}

\subsection{Calibration Methodology}

\begin{itemize}[nosep]
    \item \textbf{Rabin (2000)} - ``Risk Aversion and Expected-Utility Theory: A Calibration Theorem'' \citep{rabin2000risk} - \textit{Foundational calibration methodology}
    \item \textbf{Camerer (2012)} - ``The Calibration of Expert Judgment'' \citep{camerer2012expertise} - \textit{Expert calibration principles}
    \item \textbf{List \& Shogren (2000)} - ``Calibration of Willingness-to-Pay'' \citep{list2000revealed} - \textit{Field calibration methods}
\end{itemize}

\subsection{Key Parameter Sources (Tier 1)}

\begin{itemize}[nosep]
    \item \textbf{Loss Aversion ($\lambda$):}
    \begin{itemize}[nosep]
        \item Kahneman \& Tversky (1991) - $\lambda \approx 2.25$ \citep{kahneman1991loss}
        \item Benartzi \& Thaler (1995) - Myopic loss aversion \citep{benartzi1995myopic}
    \end{itemize}
    \item \textbf{Present Bias ($\beta$):}
    \begin{itemize}[nosep]
        \item Laibson (1997) - Golden eggs model \citep{laibson1997golden}
        \item Ericson \& Laibson (2017) - Intertemporal choice handbook \citep{ericson2017money}
    \end{itemize}
    \item \textbf{Context Effects:}
    \begin{itemize}[nosep]
        \item Choi et al. (2004) - Default effects \citep{choi2004optimal}
        \item Gabaix \& Laibson (2006) - Shrouded attributes \citep{gabaix2006shrouded}
    \end{itemize}
\end{itemize}

\subsection{Application and Validation}

\begin{itemize}[nosep]
    \item \textbf{DellaVigna \& Linos (2022)} - ``RCTs to Scale'' \citep{dellavigna2022rcts} - \textit{Large-scale nudge validation}
    \item \textbf{DellaVigna \& Pope (2018)} - ``What Motivates Effort?'' \citep{dellavigna2018motivates} - \textit{Expert forecasts vs. actual effects}
\end{itemize}

\medskip
\textbf{Complete Parameter Bibliography:} See Appendix UNMAPPED_EST for the full 200+ parameter values with sources.

% =============================================================================
% CROSS-REFERENCE FOOTER
% =============================================================================
\vfill
\begin{tcolorbox}[colback=gray!10!white, colframe=gray!50!black]
\small
\textbf{Cross-References:}
BBB/EST (Parameters) $\bullet$
CAL (Calibration) $\bullet$
DP (Pipeline) $\bullet$
FM (Formalization) $\bullet$
DR (Data Requirements) $\bullet$
G (Glossary)
\end{tcolorbox}

\end{chapter}
